\icode{E-O-SI-L3-261003-002}
为了研究不同材料的透光性，实验小组将折射率为\(n\)的材料做成厚度极小、半径为\(R\)的圆置于空气中。现将一点光源置于该材料内部，会发现光源处于某些位置时发出的光第一次到达圆边界时全部能够折射出圆边界，实验小组将能满足这样现象的光源位置组成的区域命名为“全亮区”，则全亮区的面积是\paren[]

\begin{tikzpicture}[font=\scriptsize]
	\draw (0,0) circle(1.5);
	\fill (0,0) circle(.05);
	\node at(0,-1){\(n\)};
	\node at(1.5,-1.2){空气};
\end{tikzpicture}
\begin{choices}
	\item \(\frac{\pi R^2}{n^2}\)
	\item \(\frac{\pi R^2}{4n^2}\)
	\item \(\pi n^2R^2\)
	\item \(\frac{\sqrt{2}}{2}\pi n^2R^2 \)
\end{choices}